\section{问题分析}

% 参考优秀论文 A196：问题分析按小问逐项展开。每问通常使用逻辑连续的自然段，
% 形成“当前对象/条件 → 真正难点 → 数学抓手 → 建模转化 → 准备建立的结构
% → 与前问关系”的承接，不写成散乱短句、名词清单或软件流水线；
% 不提前展开完整公式、算法细节和最终结果。

\subsection{问题一的分析}
% 说明问题一的研究对象、关键困难、所需判定/估计/优化量，以及为什么需要后续模型准备或本问模型。

\subsection{问题二的分析}
% 若依赖问题一，明确继承了什么结果、还新增什么变量/目标/约束；不要重新复述问题一全部机制。

% 按实际小问数量复制，不固定五问：
% \subsection{问题三的分析}
% \subsection{问题四的分析}
% \subsection{问题五的分析}

% 本章不写“先读取数据、再建模、再求解、再绘图”的软件流水线；
% 应从题目对象和数学困难出发形成各问之间的递进关系。
